\documentclass[11pt]{article}

\usepackage[margin=1in]{geometry}
\usepackage{array}
\usepackage{booktabs}
\usepackage{enumitem}
\usepackage{hyperref}
\usepackage{listings}
\usepackage{microtype}
\usepackage{needspace}
\usepackage{tabularx}
\usepackage[T1]{fontenc}
\usepackage{lmodern}
\usepackage{xcolor}

\hypersetup{
    colorlinks=true,
    linkcolor=blue,
    urlcolor=blue
}

\definecolor{codebg}{RGB}{247,247,247}

\newcolumntype{Y}{>{\raggedright\arraybackslash}X}
\newcolumntype{L}[1]{>{\raggedright\arraybackslash}p{#1}}

\lstset{
    basicstyle=\ttfamily\small,
    backgroundcolor=\color{codebg},
    frame=single,
    framerule=0.3pt,
    rulecolor=\color{black!25},
    columns=fullflexible,
    keepspaces=true,
    showstringspaces=false,
    breaklines=true,
    xleftmargin=0.5em,
    xrightmargin=0.5em,
    aboveskip=0.8em,
    belowskip=0.8em
}

% A conjecture to challenge (Rust-domain principle 10): the claim, the
% observation that would falsify it, and why falsification would matter.
\newenvironment{claims}{\begin{description}[leftmargin=1em,style=nextline,itemsep=0.4em]}{\end{description}}
\newcommand{\claim}[4]{\item[Claim (#1).] #2\\
    \textit{Falsified by:} #3\\
    \textit{Matters because:} #4}

\setlist[itemize]{itemsep=0.15em,topsep=0.3em}

\title{\textit{cs-mail} Product Build Plan\\
\large Server, Administration, Member Desktop, and Sender Request Page}
\author{C-SQD}
\date{Draft 0.5 --- 25 September 2026}

\setlength{\emergencystretch}{1.5em}

\begin{document}
\maketitle

\begin{abstract}
This plan turns the implemented Rust domain into a runnable product. It builds
the core relationship-request loop end to end first---operator service, identity
enrollment, request, decision, and settlement, driven from a command-line
client---and then adds depth: a public request page for senders without an
account, the macOS member application, recovery across time and failure, the
annual service and member-cycle experience, and a small pilot. cs-mail's core is
also required by the future Phoros Account, so user testing informs design rather
than deciding whether development continues. The pilot introduces no real
redeemable value.
\end{abstract}

\medskip\noindent\fbox{\parbox{\dimexpr\linewidth-2\fboxsep-2\fboxrule\relax}{\small\textbf{Design principles.} Before changing the implementation or this design, read the \href{docs/rust-domain-principles.md}{Rust-domain design principles} (\texttt{docs/rust-domain-principles.md}). They are binding for cs-mail and identity-model; new work must not regress them.}}\medskip

\section{What Changed From Draft 0.4}

Draft~0.4 built complete layers in sequence and reached the first relationship
request only at its third and fourth milestones. This draft reorders the work:

\begin{itemize}
    \item The core request loop runs end to end from a developer CLI at
    Milestone~2, before any graphical client.
    \item Enrollment and annual service use the existing identity-model and
    billing domains from the start, with simulated evidence supplied by
    development support code. No fixture or bypass path is later replaced.
    \item A public \emph{sender request page} lets a person without a cs-mail
    account pick a request class, pay, and write a message. It first requires a
    guest-sender domain design (Milestone~3a).
    \item The macOS application follows the proven loop. Billing and distribution
    \emph{experience} moves later; its domain rules are already implemented.
    \item Each milestone names the conjectures its tests challenge, following
    principle~10 of the \href{docs/rust-domain-principles.md}{Rust-domain
    principles}, plus a learning question that does not gate progress.
    \item Development and testing are local. There is no hosted CI.
\end{itemize}

\section{Build on the Implemented Domain Contracts}

The \textit{Protocol Specification} governs requests and permission. The normative
C-SQD financial-program section of the \textit{Deployment and Migration Profile}
governs annual utility billing, the pool, and equal member distributions. The
\textit{Rust Reference Architecture} describes the current ownership boundaries.
Start with the \href{DOMAIN_MODEL.md}{C-SQD domain model} for confirmed rules and
deferred choices, the \href{docs/product-identity-model.md}{product identity
model} for accounts and logins, and the
\href{docs/rust-domain-principles.md}{Rust-domain principles} for how new code is
shaped and validated.

The workspace already implements the request lifecycle with recipient-defined
request classes, lane acceptance with atomic refund settlement, native
correspondence, consent-scoped decryption, request-specific payment operations
and refunds, forfeiture maturity, annual allocation, fixed service contracts with
advance collection, funded coverage, and automatic lump-sum distribution
preparation. Signed service and client APIs, PostgreSQL transactions, and worker
libraries exercise those rules with a simulated external provider. They are not
yet a running server or a user-facing application.

Product work reuses those owners and APIs. It does not recreate deprecated labels,
aliases, fallback readers, or parallel implementations. When product work shows
that a domain rule is wrong or incomplete, the change is made as an explicit,
reviewed cutover in the owning domain, not patched in an adapter.

\section{The Product and the Pilot}

\paragraph{Clients and their roles.}
\begin{itemize}
    \item \textbf{Member application (macOS, Tauri).} Holds endpoint keys, the
    encrypted inbox, request decisions, request classes, lanes, blocks, and
    financial status. Native content is end-to-end encrypted.
    \item \textbf{Sender request page (web).} Lets a person without a cs-mail
    account select a recipient's published request class, authorize one
    conditional charge, and write a message. The page is served by the provider,
    so it cannot carry the end-to-end promise; its messages are labeled as an
    explicit privacy downgrade, like inbound SMTP.
    \item \textbf{Developer CLI.} Drives member journeys before the desktop
    application exists and remains a development and operations tool.
    \item \textbf{Administration CLI.} An authenticated operator client.
\end{itemize}

A full web inbox for members is not planned. It would place endpoint keys in code
delivered by the provider on each page load, which conflicts with the promise that
the provider cannot decrypt content. The client core compiles to WebAssembly, so
the option remains available if that trust model is ever accepted or mitigated.

\paragraph{Pilot.} An invite-only group served by one provider, one region, and
one PostgreSQL deployment. Recipients are recruited from people who receive
substantial unsolicited contact; senders are people trying to reach them,
including guests once the request page ships.

\paragraph{Value.} A deterministic external-payment simulator approves,
declines, delays, or reports an uncertain charge; confirms or fails a refund; and
executes or fails a member payment. All amounts are explicitly nonredeemable.
Users receive no spendable test-credit wallet; the operator configures simulator
scenarios.

\paragraph{Out of scope for the pilot.} Independent-provider federation, a full
SMTP relay, public signup, live payments, mobile releases, and multi-region
operation.

\section{Architecture of the Product Layer}

\begin{lstlisting}
crates/
  existing cs-mail-* crates     # domain, application, storage, workers
  cs-mail-client                # extended client core (ports for keys,
                                #   local storage, transport, clock)
  cs-mail-api-types             # add when transport contracts need it
  cs-mail-mobile-ffi            # Milestone 8
apps/
  cs-maild                      # HTTP API, workers, sender page hosting
  cs-mail-admin                 # authenticated operator CLI
  cs-mail-cli                   # developer member client
  cs-mail-desktop               # Tauri shell and macOS adapters
  cs-mail-web                   # sender request page (static assets)
ui/                             # interface components shared by
                                #   desktop and web
../identity-model               # separate repository, sibling checkout
\end{lstlisting}

Dependencies point inward: applications use the client core and application
layer, which use the domain; domain crates never depend on applications.
\textbf{Applications are adapters.} They validate input, invoke application use
cases, and render outcomes. They make no domain decisions, obtain trusted time
only as the application-persistence boundary prescribes, and never translate a
signed command into a different meaning. The client core compiles to WebAssembly;
the member application nonetheless runs it natively through Tauri.

identity-model remains a separate repository because it will also serve the
Phoros Account. cs-mail pins the reviewed identity sources in
\texttt{identity-source.sha256}; commit identity-model changes before the cs-mail
changes that depend on them and keep the manifest in step.

\section{Contracts That Apply to Every Milestone}

\paragraph{Keys and authority.} Endpoint private keys use operating-system secure
storage (or, for the developer CLI only, a local-file implementation of the same
key-custody port) and never enter logs, configuration, SQLite, provider APIs, or
diagnostics. Device enrollment and revocation require authenticated authority.
Administrative credentials are distinct from user credentials; the admin CLI uses
an authenticated API, never direct table writes.

\paragraph{Evidence and simulation.} Simulated login, bank-ownership, funding, and
payment evidence, together with scenario builders, live in development support
code. They enter the product through the same ports as real evidence and are never
presented as business commitments. Rates, thresholds, prices, and dates use
explicitly versioned simulation policies.

\paragraph{Outcomes and uncertainty.} Every visible outcome has a durable cause
and a stable retry identity. ``Outcome unknown,'' ``failed,'' ``expired,'' and
``denied'' remain distinct through storage, API, and interface. The interface
separates submission preparation, awaiting a decision, refund due, refund
completed, annual service coverage, finalized member allocation, and payment
completion. A distribution never marks a service year unpaid. No screen suggests
that rejecting a sender pays the recipient or that a payable can fund a request.

\paragraph{Versions and data.} Protocol, wire, API, event, policy, and database
versions are separate. Until the fresh-database policy ends (decided at
Milestone~0), schema changes may replace earlier migrations; afterwards
migrations are forward-only. Logs contain scoped references and typed errors, not
plaintext, secrets, identity mappings, full signed commands, or complete financial
snapshots. Retention policy is explicit before pilot data exists.

\paragraph{Validation.} Tests challenge named conjectures whose falsification
would matter; each milestone lists its claims and agrees its testing scope before
work begins. Claims about locking, atomicity, or concurrency are challenged
against PostgreSQL, not an in-memory adapter. Guarantees already enforced by types
need no redundant runtime test. Demonstration scripts show a journey; they are not
evidence for a claim. At each handoff, run the local quality gate:
\begin{lstlisting}
cargo fmt --all -- --check
cargo clippy --workspace --all-targets -- -D warnings
cargo test --workspace
CS_MAIL_TEST_DATABASE_URL=... cargo test --workspace -- --ignored
\end{lstlisting}

\section{Milestone Map}

\begin{center}
\small
\begin{tabularx}{\textwidth}{@{}L{0.07\textwidth}L{0.33\textwidth}Y@{}}
\toprule
Stage & Result & Exit evidence \\
\midrule
0 & Product skeleton & Buildable apps, inward dependencies, quality gate \\
1 & Service and identity boot & Clean start, restart, invites, simulator scenarios \\
2 & Core loop from the CLI & Every request outcome, from a clean database \\
3a & Guest-sender domain design & Agreed types, evidence, ownership, and rules \\
3b & Sender request page & A guest completes a request; every outcome settles \\
4 & macOS member application & Milestone 2 journeys in the app, keys in keychain \\
5 & Time and failure & Crash, race, and deadline claims challenged \\
6 & Annual service and member cycle & A full simulated year, including nonrenewal \\
7 & Pilot & Signed builds, recovery, restore, pilot findings \\
8 & Portability and Phoros & Windows, Swift/Kotlin parity, Phoros integration \\
\bottomrule
\end{tabularx}
\end{center}

Each stage depends on the preceding one, except that Milestone~4 may proceed
before 3b if the guest-sender design remains unresolved. A milestone is complete
when its demonstration runs from a clean checkout and its claims have been
challenged, not merely when new code compiles.

\section{Milestone 0: Product Skeleton}

Create \texttt{apps/cs-maild}, \texttt{apps/cs-mail-admin}, and
\texttt{apps/cs-mail-cli} with shared package metadata, lints, and dependency
versions. Each reports its version and help. Add a local script that runs the
full quality gate without hiding individual checks, and document the inward
dependency rule.

\textbf{Decision to record:} when the fresh-database policy ends and migrations
become forward-only---at the latest, before any pilot data is created.

\textbf{Exit:} clean checkout, passing gate, running binaries. No listener,
enrollment, or domain change belongs to this stage.

\section{Milestone 1: The Service and Identity Boot}

The operator can migrate an empty database, bootstrap one provider idempotently,
start the API, create invites, publish service offers and request pricing, and
select simulator scenarios---without editing SQL. Use Tokio, Axum, Rustls,
structured tracing, and Clap. Keep synchronous storage behind a bounded blocking
boundary until a tested replacement is warranted. Readiness distinguishes process
health from database readiness.

The identity-model enrollment service runs beside \texttt{cs-maild}, with a
development OIDC provider (the Keycloak realm in \texttt{identity-model/ci}) and
simulated bank-ownership evidence from development support code. The payment
simulator persists operation identities and supports delayed, duplicate,
out-of-order, and failed results; its scenario changes are audited and separate
from canonical financial writes.

\begin{claims}
\claim{PostgreSQL}{Repeating bootstrap cannot create a second provider authority or replace the key registry.}
    {Two bootstrap runs, including one interrupted midway, leave two authorities or a changed registry.}
    {Duplicate authority would let two keys speak for the provider.}
\claim{PostgreSQL}{Replaying or concurrently redeeming one invite cannot enroll two accounts.}
    {Two concurrent redemptions both succeed.}
    {One account per person is an anti-abuse objective and a member-share rule.}
\claim{PostgreSQL}{A restart cannot duplicate a simulated payment operation.}
    {A killed and restarted service shows two operations for one operation identity.}
    {Duplicate operations become duplicate charges once payments are real.}
\end{claims}

\textbf{Exit:} a clean-start transcript, restart demonstration, and redacted
diagnostics.

\section{Milestone 2: The Core Loop From the Command Line}

Alice and Bob enroll through identity-model using the developer CLI and establish
their protocol identities and device keys. Each purchases annual service through
the existing billing domain; the simulator finalizes funding and coverage begins
within the purchased interval.

Bob publishes request classes. Alice resolves Bob, reviews $C$, $S$, the total
conditional charge, the decision deadline, and possible refunds; authorizes one
charge; and submits an encrypted initial message after funding reaches finality.
Bob accepts with standing permission or an express lane, rejects, blocks, or lets
the request expire. Alice sends follow-ups when allowed. Both sides see settlement:
refund owed and confirmed, processing revenue retained, collateral pending
forfeiture. Uncertain outcomes are shown as uncertain.

\begin{claims}
\claim{PostgreSQL}{Retrying a submission whose outcome is unknown cannot create a second request or charge.}
    {A client that loses the response and retries produces two requests or two captures.}
    {One request must remain one charge.}
\claim{PostgreSQL}{Acceptance commits permission and the full refund obligation together or not at all.}
    {A failure injected during acceptance leaves a lane or permission without its refund obligation, or the reverse.}
    {Either half alone breaks the settlement contract.}
\claim{in-memory and PostgreSQL}{An altered retry under an existing operation identity is rejected.}
    {A retry with changed terms or content under the same identity is accepted.}
    {Retry identities would stop protecting request terms.}
\end{claims}

\textbf{Exit:} a scripted two-member demonstration of every outcome from a clean
database. \textbf{Learning:} does the decision flow read clearly, and which wording
confuses?

\section{Milestone 3a: Guest-Sender Domain Design}

The current domain recognizes two senders: members, with an account, verified bank
association, and protocol identity; and organizational legacy senders, identified
by a DMARC-authenticated domain. A person using the request page without an
account is neither. Before any request-page code, record in the domain model and a
design note:

\begin{itemize}
    \item the guest sender as a distinct type, and what evidence (for example a
    verified email address or payment evidence) identifies one and what that
    evidence establishes;
    \item who owns guest eligibility and anti-abuse history, given that
    principal-recipient history is keyed on a principal a guest does not have;
    \item what acceptance grants a guest---permission bound to verified
    evidence, a path to membership, or a bounded lane;
    \item how refunds reach a guest's payment method without an account;
    \item how guest messages are marked as non-end-to-end and how that label
    reaches the recipient.
\end{itemize}

Follow the existing treatment of legacy senders where it fits. This is the one
planned change to the core domain; it is made as a reviewed cutover, not an
adapter-level extension. \textbf{Exit:} the agreed design, with its conjectures
stated. If it remains unresolved, Milestone~4 proceeds first.

\section{Milestone 3b: Sender Request Page}

Each member has a public request link. A guest opens it, selects a published
class, reviews $C$, $S$, the refund rules and the privacy label, authorizes one
simulated charge, and writes a message. Encryption to the recipient's key on
arrival is an application use case; plaintext is not stored or logged. The
recipient decides from the developer CLI until Milestone~4. The guest receives
outcome and refund notices by email.

Optionally, unknown senders emailing a member's address receive an automatic
reply with the request link, using \texttt{cs-mail-smtp-gateway} for sender
authentication. This is in scope only if relay operations stay small.

\begin{claims}
\claim{PostgreSQL}{A guest cannot exceed pending-request limits for one recipient by changing email addresses.}
    {Repeated requests under fresh addresses from the same payment evidence all reach the recipient. The strength of this claim is set by 3a.}
    {The request page would become a spam channel.}
\claim{inspection and integration}{Guest-message plaintext is never written to storage, logs, or diagnostics.}
    {A secret-scan of the database, logs, and diagnostics after a guest request finds the message text.}
    {The downgrade label would understate the exposure.}
\end{claims}

\textbf{Exit:} a guest completes a request end to end and every outcome settles.
\textbf{Learning:} do strangers understand the page, and what collateral feels
reasonable?

\section{Milestone 4: macOS Member Application}

Alice installs the application, enrolls or recovers through identity-model, and
reopens the same account after quitting. Keys live in the macOS keychain behind
the key-custody port; SQLite holds ciphertext, key references, cursors, minimized
relationship metadata, and nonsecret preferences only. Portable client use cases
coordinate enrollment and device management; the Tauri shell stays thin.

Bob sees the chosen class, one request with its grouped messages, and a clear
deadline, and decides with permission explained first and money second. He
manages request classes with collateral from the operator's bounded menu,
follow-up preferences, lanes, and blocks. Financial status shows request charges
and refunds separately from annual service and distributions. Interface
components are shared with the request page.

\begin{claims}
\claim{integration}{Secret key material never reaches SQLite, logs, or diagnostics.}
    {A secret-scan after enrollment, messaging, and export finds key bytes.}
    {The endpoint-encryption promise would be void.}
\claim{integration}{After quit, reopen, and projection rebuild, every displayed outcome matches the server's canonical record.}
    {A decision, refund, or uncertain charge appears differently after reopening.}
    {Users would act on false financial or permission state.}
\end{claims}

\textbf{Exit:} Milestone~2's journeys in the application, including keychain
denial and quit/reopen. \textbf{Learning:} do recipients publish classes, how
many, and how long do decisions take?

\section{Milestone 5: Recover Across Time and Failure}

Run scheduled deadlines, outbox delivery, payment reconciliation, and retention
through supervised leased workers with a distinct scheduler authority, bounded
retries, graceful shutdown, database TLS, and a bounded connection pool. Jobs
recheck current state and use stable operation identities. Offline recipients and
cursor replay must not duplicate delivery or financial effects.

\begin{claims}
\claim{PostgreSQL, crash injection}{A crash between capture and recording its response cannot lose or duplicate the charge.}
    {After restart, reconciliation shows no capture or two captures for one request.}
    {Money would be lost or taken twice.}
\claim{PostgreSQL, race}{Cancellation racing a delayed funding confirmation returns captured money exactly once.}
    {The request is both cancelled and submitted, or the refund is issued twice or never.}
    {Cancellation is the user's safety valve before submission.}
\claim{PostgreSQL}{Out-of-order provider events cannot reopen a terminal settlement.}
    {A late event changes a settled request's obligations.}
    {Terminal financial outcomes must stay fixed.}
\claim{PostgreSQL}{Coverage expiring while a request is unsettled stops new communication but not decision or settlement.}
    {The recipient cannot decide, or a refund stalls, after coverage ends.}
    {Obligations outlive service coverage.}
\end{claims}

\textbf{Decision to record:} the bounded-acceptance timing edge cases currently
deferred in the domain model.

\section{Milestone 6: Annual Service and Member-Cycle Experience}

Build the billing and distribution experience and run collection retries in
operation. A member reviews a published offer with annual price, advance
collection date, and fixed service interval; failed and late collection retries
preserve the price and interval; a pending charge does not imply renewed service.

Finalize a simulated UTC calendar year from fixed funding and eligibility inputs,
showing the once-only corporate assessment, equal minor-unit allocations, and
restricted remainder. Pay each member one lump sum of the previous year's share
to the associated verified bank account. With $U=\$120$, allocation scenarios of
\$80, \$120, and \$170 each produce one payment, reported as rebate/excess of
\$80/\$0, \$120/\$0, and \$120/\$50. Repeat with nonrenewal and closure,
preserving the minimal payment route.

\begin{claims}
\claim{PostgreSQL}{A failed or unanswered payout attempt cannot create a second allocation or payable.}
    {After failure and restart, a member has two payables for one pool year.}
    {Pool funds would be distributed twice.}
\claim{PostgreSQL}{A later allocation year cannot reassess or re-divide amounts already assigned.}
    {Year $n+1$ changes an assigned year-$n$ amount.}
    {Members' fixed shares would be unstable.}
\claim{integration}{A distribution never marks a paid service year unpaid or revokes coverage.}
    {Coverage or charge status changes when a distribution is paid.}
    {Service and distribution are separate transactions by rule.}
\end{claims}

\textbf{Exit:} a full simulated year, including a nonrenewing member who is still
paid.

\section{Milestone 7: Run a Small Supported Pilot}

Package and sign the macOS application and versioned server and CLI artifacts.
Complete user-held recovery and trusted-device enrollment with explicit device
revocation and no provider decryption authority. Provide onboarding for annual
service, requests, refunds, simulated distributions, follow-up preferences,
encryption, recovery, and device loss; scoped diagnostics export; support and
abuse intake; deletion procedures; backup and restore; and release rollback.
Backups and retained obligation records respect separate keys and access roles.

Recruit recipients who receive substantial unsolicited contact and senders who
want to reach them, including guests if Milestone~3b has shipped. Nondeveloper
testers use clean machines and the pilot guide without direct database repair.

\begin{claims}
\claim{integration}{Restore followed by reconciliation reproduces every contract, request, refund, allocation, payable, and inbox state.}
    {Any obligation differs between the pre-backup and reconciled post-restore state.}
    {Restore is the operator's last line of defense.}
\claim{integration}{A revoked device cannot authorize commands or decrypt new content.}
    {A command signed by a revoked device is accepted, or new content is readable to it.}
    {Device loss is a routine event.}
\end{claims}

\textbf{Exit:} signed releases, restore and recovery records, resolved critical
usability findings, and a written summary of pilot findings.

\section{Milestone 8: Portability and Phoros}

Deliver Windows adapters once the macOS path is stable. Expose a narrow FFI facade
for identity, requests, decisions, encrypted inbox, billing and coverage,
financial status, distributions, and devices, with bounded values, stable error
codes, cancellation, and explicit ownership. Generate Swift and Kotlin bindings
with mock platform adapters and run the same client-core stories through each.
Add a live payment connector.

For Phoros, its enrollment ceremony adopts an existing cs-mail subject through
verified account control, as the product identity model prescribes. Version the
identity contract as a published interface once two products consume it. Design
cross-product account linking on explicit, granular, revocable consent.

\section{Deferred and Open}

\begin{itemize}
    \item Bank-account replacement, full account management, and detailed
    erasure policy.
    \item Independent-provider federation and full SMTP relay operations.
    \item A full web inbox for members.
\end{itemize}

\section{How an Implementation Agent Should Proceed}

For each slice, read the applicable contracts, the Rust-domain principles, and
local instructions. Name one user-visible outcome, identify its trust and
persistence boundaries, agree the testing scope, then implement contracts,
application orchestration, adapters, and interface in that order. Report what
changed, what was demonstrated, which claims were challenged, and what remains
open or untested.

Do not bypass signed execution, write canonical tables from a CLI, persist
plaintext for convenience, hide an uncertain charge, add tests merely to
accompany code, or weaken an existing check to finish a milestone. A domain change
requires an explicit decision and a reviewed cutover. Scope does not expand to
live value, federation, or a full SMTP relay because an adapter exists.

If reconciliation fails, preserve the affected evidence and investigate the
specific obligation. If secrets leak, stop accumulating affected data until the
boundary is repaired. If a policy parameter is unresolved, use an explicitly
versioned simulation policy and keep the live deployment gate closed.

\section{First Assignment}

Milestone~0, followed by the smallest Milestone~1 path: start the service,
establish authenticated administration, run identity enrollment against the
development OIDC provider, and expose one durable simulated annual charge. Later
milestones extend that path without reimplementing domain rules.

\end{document}
